\documentclass[11pt]{amsart}

\usepackage[margin=1.15in]{geometry}
\usepackage{amsmath,amssymb,amsthm,mathtools}
\usepackage{microtype}
\usepackage{xurl}
\usepackage{xcolor}
\usepackage{listings}
\usepackage[hidelinks]{hyperref}
\usepackage{aliascnt}
\usepackage[nameinlink,noabbrev]{cleveref}

\definecolor{keywordcolor}{RGB}{88,41,120}
\definecolor{sortcolor}{RGB}{35,79,140}
\definecolor{symbolcolor}{RGB}{125,48,48}
\definecolor{commentcolor}{RGB}{42,105,62}
% Compact Lean language for the listings package, adapted from the
% user-supplied lstlean.tex definition.
\lstdefinelanguage{lean}{
  morekeywords={import,namespace,section,end,open,variable,variables,
    theorem,lemma,corollary,example,def,definition,structure,class,
    instance,inductive,where,by,have,show,exact,apply,refine,
    let,in,if,then,else,match,with,fun,forall,exists,noncomputable,
    protected,private,local,attribute,set_option,universe,universes,
    check},
  morekeywords=[2]{Type,Prop},
  sensitive=true,
  morecomment=[l][\itshape\color{commentcolor}]{--},
  morecomment=[s][\color{commentcolor}]{/-}{-/},
  morestring=[b]",
  keywordstyle=\color{keywordcolor},
  keywordstyle=[2]\color{sortcolor},
  identifierstyle=\color{black},
  stringstyle=\color{symbolcolor},
  showstringspaces=false,
  keepspaces=true,
  breaklines=true,
  columns=fullflexible,
  literate=
    {¬}{{$\neg$}}1
    {≠}{{$\neq$}}1
    {·}{{$\cdot$}}1
    {ι}{{$\iota$}}1
    {⁺}{{$^{+}$}}1
    {₀}{{$_0$}}1
    {₁}{{$_1$}}1
    {₂}{{$_2$}}1
    {₃}{{$_3$}}1
    {ℕ}{{$\mathbb{N}$}}1
    {ℝ}{{$\mathbb{R}$}}1
    {ϖ}{{$\varpi$}}1
    {→}{{$\to$}}1
    {↔}{{$\leftrightarrow$}}1
    {↥}{{$\uparrow$}}1
    {∀}{{$\forall$}}1
    {∃}{{$\exists$}}1
    {∈}{{$\in$}}1
    {∑}{{$\sum$}}1
    {∧}{{$\wedge$}}1
    {≤}{{$\leq$}}1
    {⊆}{{$\subseteq$}}1
    {⋃}{{$\bigcup$}}1
    {⋙}{{$\mathbin{\gg}$}}1
    {𝒪}{{$\mathcal O$}}1
    {⟨}{{$\langle$}}1
    {⟩}{{$\rangle$}}1
}

\lstset{
  language=lean,
  basicstyle=\ttfamily\footnotesize,
  frame=single,
  framerule=0.25pt,
  rulecolor=\color{black!35},
  backgroundcolor=\color{black!2},
  xleftmargin=0.7em,
  xrightmargin=0.7em,
  aboveskip=0.65\baselineskip,
  belowskip=0.65\baselineskip,
  breaklines=true,
  columns=fullflexible,
  captionpos=b
}

\newtheorem{theorem}{Theorem}[section]
\newaliascnt{proposition}{theorem}
\newtheorem{proposition}[proposition]{Proposition}
\aliascntresetthe{proposition}
\newaliascnt{lemma}{theorem}
\newtheorem{lemma}[lemma]{Lemma}
\aliascntresetthe{lemma}
\newaliascnt{corollary}{theorem}
\newtheorem{corollary}[corollary]{Corollary}
\aliascntresetthe{corollary}
\theoremstyle{definition}
\newaliascnt{definition}{theorem}
\newtheorem{definition}[definition]{Definition}
\aliascntresetthe{definition}
\newaliascnt{remark}{theorem}
\newtheorem{remark}[remark]{Remark}
\aliascntresetthe{remark}
\newcommand{\Spa}{\operatorname{Spa}}
\newcommand{\im}{\operatorname{im}}
\newcommand{\cO}{\mathcal O}
\newcommand{\onto}{\twoheadrightarrow}

\title{Uniform sheafy Tate domains that are not stably uniform}
\author{Christopher Birkbeck}
\date{}

\begin{document}

\begin{abstract}
Buzzard--Verberkmoes and Mihara show that a stably uniform Tate Huber
ring is sheafy.  We construct two uniform sheafy Tate domains which are
not stably uniform, answering Question~7 of Kedlaya's
\emph{Nonarchimedean Scottish Book}.  In the first example a rational
localisation is nonreduced, whilst in the second the nonuniform rational
localisation is an integral domain.  The two main results are due to
ChatGPT 5.6 Sol.
Both constructions, including their sheafiness and the rational localisations
witnessing failure of stable uniformity, have been formalised in Lean~4.
\end{abstract}

\subjclass[2020]{Primary 14G22; Secondary 13J10}
\keywords{Huber ring, adic spectrum, sheafiness, stable uniformity, rational localisation, Milnor square}

\maketitle

\section{Introduction}

The point at issue is the behaviour of rational localisations.  The examples
of Buzzard--Verberkmoes and Mihara show that a uniform Tate Huber ring can lose
uniformity on a rational chart and, separately, that uniformity alone does not
imply sheafiness \cite{BV,Mihara}.  They leave open whether sheafiness itself
prevents this loss of uniformity.  Hansen and Kedlaya ask this explicitly
\cite[Remark~3.16]{HK}; it is also Question~7 of Kedlaya's
\emph{Nonarchimedean Scottish Book} \cite{Scottish}.  The answer is no.

ChatGPT 5.6 Sol found both examples and supplied the proofs.  The Lean
formalisation was developed with ChatGPT.  We have aimed to cite related work
that may have informed this process; see \cref{rem:earlier-work}.
\Cref{sec:discovery} gives a fuller account, and \cref{sec:formalisation}
records the formal definitions and theorem statements.

We begin with the ``finite-jet pullback''.  This is a
nonnoetherian uniform integral domain for which a distinguished rational
localisation is nonreduced, and hence is no longer uniform.  Let $k$ be a
complete discretely valued nonarchimedean field, with valuation ring
$k^\circ$, uniformizer $\varpi$, and residue field $\widetilde{k}$.  Put
\[
  L_0=k^\circ\langle W,W^{-1}\rangle,
  \qquad
  B_0=k^\circ\langle W,Q\rangle/(Q^2),
\]
\[
  C_0=L_0\langle Q\rangle,
  \qquad
  D_0=L_0\langle Q\rangle/(Q^2),
\]
where $k^\circ\langle W,W^{-1}\rangle$ means
$k^\circ\langle W,V\rangle/(WV-1)$.  There are natural maps
$B_0\to D_0$ and $C_0\onto D_0$.  Define
\[
  A_0=B_0\times_{D_0}C_0,
  \qquad
  A=A_0[1/\varpi],
\]
and endow $A$ with the topology defined by the ring of definition $A_0$.

The main result is the following.

\begin{theorem}\label{thm:main}
The ring $A$ is a complete uniform nonnoetherian Tate $k$-algebra and an
integral domain, with $A^\circ=A_0$.  The Huber pair $(A,A^\circ)$ is sheafy.
However,
\[
  A\left\langle \frac{W}{\varpi}\right\rangle
  \cong k\langle X,Q\rangle/(Q^2),
  \qquad X=\frac{W}{\varpi},
\]
so $A$ is not stably uniform.
\end{theorem}

The proof uses Milnor squares of complete Tate rings.  Recall that a cartesian square of
rings with a surjective leg is called a Milnor square \cite[Section~2]{Milnor}; the
complete Tate setting is used by Kerz--Saito--Tamme in their work on analytic
$K$-theory \cite[Section~3]{KST}.

\begin{definition}\label{def:strict-milnor}
A commutative square of complete Tate $k$-algebras
\[
\begin{matrix}
R&\longrightarrow&C\\
\downarrow&&\downarrow\\
B&\longrightarrow&D
\end{matrix}
\]
is a \emph{strict Milnor square} if it is cartesian, the sequence of underlying topological $k$-vector spaces
\[
  0\longrightarrow R\longrightarrow B\oplus C
  \xrightarrow{\ (b,c)\mapsto \bar b-\bar c\ }D\longrightarrow 0
\]
is strict exact, and $C\to D$ is a strict surjection.
\end{definition}

After rotating the square if necessary, the initial Milnor square is strict by the
nonarchimedean open mapping theorem \cite[Lemma~3.1]{KST}.  Our main input is that,
for this finite-jet square, the strict exact row persists after completed rational
localisation and is compatible with refinement.  This transfers sheafiness from
$B,C,D$ to $A$.  The second example shows that the failure of stable uniformity
need not come from a nonreduced localisation.

\begin{remark}[Relation with earlier work]\label{rem:earlier-work}
It is difficult to understand exactly what sources an AI agent such as
ChatGPT may have used in doing this work.  In the interest of giving
appropriate credit, we have listed here closely related work that may have
been used in the development of the results here.

Buzzard--Verberkmoes construct a uniform affinoid ring with a nonuniform rational
localisation and, separately, a uniform affinoid ring whose structure presheaf is not a
sheaf \cite[Propositions 17 and 18]{BV}.  Mihara gives further examples of both
phenomena \cite{Mihara}.  These examples do not give a uniform sheafy ring which is not
stably uniform.

Ben-Bassat--Kremnizer prove strict two-term resolutions
for Weierstrass and Laurent localisations, and reduce a general rational
localisation to these cases \cite[Lemmas~5.13--5.14]{BBK}.
Bambozzi--Kremnizer formulate rational localisation in terms of Koszul
complexes and prove strict Koszul regularity
\cite[Notation~4.3, Definitions~4.5--4.6, Proposition~4.9,
Lemma~4.11, and Corollary~4.13]{BK}.  The lattice estimates in
\cref{lem:koszul} are the form needed to show that rational localisation
preserves the finite-jet Milnor square.  The resulting uniform sheafy domain
which is not stably uniform appears to be new.

The weighted-parity construction of \cref{sec:second-example} is closer to the
infinite monomial-support examples of Buzzard--Verberkmoes and Mihara.  Its additional
feature is that rational localisation can be calculated in a noetherian
finite-variable subring.
\end{remark}


\section{Conventions and notation}\label{sec:conventions}

We follow Huber for Huber pairs, adic spectra, rational subsets, and rational
localisation \cite{Huber94}, and use Wedhorn's account for strong
noetherianity and sheafiness
\cite[Proposition and Definition~6.36, Section~8.1]{Wedhorn}; the
bounded-denominator language for strictness follows Buzzard--Verberkmoes
\cite{BV}, and the quotient presentation of a rational localisation is also
discussed by Hansen--Kedlaya \cite{HK}.

\begin{itemize}
\item \emph{Tate topology.}
Throughout, $E$ denotes a complete Tate $k$-algebra.  In the discretely valued setting of this paper, this means that $E$ contains a topologically nilpotent unit and admits an open bounded $k^\circ$-subalgebra $E_0$ such that
\[
  E=E_0[1/\varpi],
  \qquad
  \{\varpi^nE_0\}_{n\geq0}
\]
is a basis of neighbourhoods of zero.  Such an $E_0$ is called a \emph{ring of definition}.  We always choose it closed, hence $\varpi$-adically complete.  A subset of $E$ is bounded if it is contained in $\varpi^{-r}E_0$ for some $r\geq0$.  If $M$ is a complete topological $k$-vector space, an open bounded $k^\circ$-submodule $M_0\subset M$ will be called a lattice.

\item \emph{Uniformity and noetherianity.}
We write $E^\circ$ for the ring of power-bounded elements.  The ring $E$ is
\emph{uniform} if $E^\circ$ is bounded, and \emph{strongly noetherian} if
$E\langle Z_1,\ldots,Z_s\rangle$ is noetherian for every $s\geq0$.  A Tate ring
is \emph{stably uniform} if every rational localisation is uniform.

\item \emph{Huber pairs and sheafiness.}
A \emph{ring of integral elements} is an open subring
$E^+\subseteq E^\circ$ which is integrally closed in $E$.  We call the Huber pair
$(E,E^+)$ \emph{sheafy} when its structure presheaf is a sheaf of complete
topological rings; in particular, the map to the matching-family equaliser for a
finite rational cover is a homeomorphism.

\item \emph{Strictness.}
We use the following bounded-denominator form of strictness.  Let $u:M\to N$ be a continuous $k$-linear map and choose lattices $M_0\subset M$ and $N_0\subset N$.  Then $u$ is strict if and only if there exists $a\geq0$ such that
\begin{equation}\label{eq:strictness}
  \varpi^a\bigl(u(M)\cap N_0\bigr)\subset u(M_0).
\end{equation}
For a surjection this becomes $\varpi^aN_0\subset u(M_0)$.  For an injection it says that the lattice induced from the target is bounded in the source.  This is the criterion used by Buzzard--Verberkmoes \cite[Lemma~2]{BV}.  We use the term \emph{strict Milnor square} in the sense of \cref{def:strict-milnor}; the underlying algebraic terminology is standard \cite{Milnor,KST}.

\item \emph{Rational localisation.}
We take every rational datum to have at least one numerator; an empty list may be
padded by $f_1=0$.  Let
$\alpha=(f_1,\dots,f_m;g)$ be a rational datum in $E$, so the ideal generated by
$g,f_1,\dots,f_m$ is open.  Since $E$ is a Tate $k$-algebra, this ideal is the unit ideal.  After multiplying the entire datum by a common power of $\varpi$, we may assume that all entries lie in $E_0$.

The rational localisation $E_\alpha$ is the separated $\varpi$-adic completion of $E[1/g]$ for the ring of definition generated by
\[
  E_0\left[\frac{f_1}{g},\dots,\frac{f_m}{g}\right].
\]
It is independent of the ring of definition and of the presentation of the rational
domain, and it is transitive under rational refinement.  Put
\[
  T_E=E\langle T_1,\dots,T_m\rangle
  =E_0\langle T_1,\dots,T_m\rangle[1/\varpi],
  \qquad r_i=gT_i-f_i,
\]
and
\begin{equation}\label{eq:graph-ideal-definition}
  I_{E,\alpha}:=(r_1,\dots,r_m)
  =\operatorname{im}\left(
    T_E^m\xrightarrow{\ d_{1,E}\ }T_E
  \right)\subset T_E,
  \qquad
  d_{1,E}(u_1,\dots,u_m)=\sum_i u_i(gT_i-f_i),
\end{equation}
Then
\begin{equation}\label{eq:graph-model}
  E_\alpha\cong T_E/\overline{I_{E,\alpha}}.
\end{equation}
Here the bar is necessary until the ideal has been shown to be closed.  This is the
standard quotient presentation of rational localisation; see
\cite[(1.2), Proposition~1.3, Lemma~1.5]{Huber94} and
\cite[Definition~3.6]{HK}.
\end{itemize}


\section{The finite-jet ring}\label{sec:the-global-ring}

In this section we identify the finite-jet algebra with a concrete subring of
$C$ and prove that it is a complete uniform domain.  Recall that
\[
  L_0=k^\circ\langle W,W^{-1}\rangle,
  \qquad
  B_0=k^\circ\langle W,Q\rangle/(Q^2),
\]
\[
  C_0=L_0\langle Q\rangle,
  \qquad
  D_0=L_0\langle Q\rangle/(Q^2),
\]
where $k^\circ\langle W,W^{-1}\rangle$ denotes
$k^\circ\langle W,V\rangle/(WV-1)$.  The natural maps
$B_0\to D_0$ and $C_0\onto D_0$ define
\[
  A_0=B_0\times_{D_0}C_0.
\]
We write
\[
  A=A_0[1/\varpi],\quad
  L=L_0[1/\varpi],\quad
  B=B_0[1/\varpi],\quad
  C=C_0[1/\varpi],\quad
  D=D_0[1/\varpi].
\]
Each of these Tate rings carries the topology defined by its subscript-zero
ring.

The map $C_0\to D_0$ has a continuous $k^\circ$-linear section given by truncation in the $Q$-variable:
\[
  f_0+Qf_1\longmapsto f_0+Qf_1.
\]
Consequently
\begin{equation}\label{eq:integral-milnor}
  0\longrightarrow A_0\longrightarrow B_0\oplus C_0
  \longrightarrow D_0\longrightarrow 0
\end{equation}
is exact, and $A_0$ is $\varpi$-adically complete.  After inverting $\varpi$ we obtain a strict Milnor square
\[
\begin{matrix}
A&\longrightarrow&C\\
\downarrow&&\downarrow\\
B&\longrightarrow&D.
\end{matrix}
\]
In fact, the chosen rings of definition already make the integral row exact, so no powers of $\varpi$ are lost in the defining square.


The map $B_0\to D_0$ is injective.  Projection to $C_0$ therefore identifies $A_0$ with the subring
\begin{equation}\label{eq:coefficient-A0}
  A_0=
  \left\{
  f_0(W)+Qf_1(W)+Q^2h(W,W^{-1},Q):
  f_0,f_1\in k^\circ\langle W\rangle,
  \ h\in C_0
  \right\}.
\end{equation}
Equivalently, $A_0$ consists of the restricted series
\begin{equation}\label{eq:monomial-support}
  \sum_{(a,b)\in S}c_{a,b}W^aQ^b,
  \qquad
  S=\{(a,b)\in\mathbb Z\times\mathbb N: b\leq 1\Rightarrow a\geq 0\},
\end{equation}
with $c_{a,b}\in k^\circ$.  Here restricted means that, for every $n$, only finitely many coefficients are nonzero modulo $\varpi^n$.  The set $S$ is closed under addition, so this is a subring of $C_0$.  The same description with coefficients in $k$ gives $A\subset C$.


Equip $A,B,C,D$ with the coefficientwise Gauss norms whose unit balls are the
displayed rings of definition.  Write
\[
  \iota_B:A\to B,\qquad \iota_C:A\to C,\qquad
  \rho_B:B\to D,\qquad \rho_C:C\to D
\]
for the structural maps, and give \(A\) the restricted Gauss norm from \(C\).  For
\(a=f_0+Qf_1+Q^2h\), the coefficient formulas give
\begin{equation}\label{eq:structural-map-bounds}
  \lVert\iota_C(a)\rVert_C=\lVert a\rVert_A,
  \qquad
  \lVert\iota_B(a)\rVert_B\leq \lVert a\rVert_A,
  \qquad
  \lVert\rho_B(b)\rVert_D\leq\lVert b\rVert_B,
  \qquad
  \lVert\rho_C(c)\rVert_D\leq\lVert c\rVert_C.
\end{equation}
Thus every structural map is norm-nonincreasing.

\begin{proposition}\label{prop:uniform-domain}
The ring $A$ is a complete uniform nonnoetherian Tate $k$-algebra and an
integral domain.  Moreover,
\[
  A^\circ=A_0.
\]
\end{proposition}

\begin{proof}
The Gauss norm on $L=k\langle W,W^{-1}\rangle$ is multiplicative, and the
restricted Gauss norm on $C=L\langle Q\rangle$ is therefore multiplicative as
well.  In particular, $C$ is a domain, and hence so is its subring $A$.  The
conditions on the constant and linear $Q$-coefficients which define $A$ are
closed conditions in $C$, so $A$ is complete.  The image of $\varpi$ in $A$
is a topologically nilpotent unit, and therefore $A$ is a Tate ring.

The coefficient description of $A_0$ shows that it is exactly the unit ball
for the restricted Gauss norm on $A$.  Since this norm is multiplicative,
\[
  \lVert a^n\rVert=\lVert a\rVert^n.
\]
It follows that $a$ is power-bounded if and only if $\lVert a\rVert\leq1$.
Thus $A^\circ=A_0$, and the unit ball is bounded, so $A$ is uniform.

For nonnoetherianity, let $J=\ker(\iota_B)=Q^2C\subset A$.  If $J$ were
generated by $a_1,\ldots,a_r$, then, for every $\ell\in L$, comparison of the
$Q^2$-coefficient in an expression
\[
  Q^2\ell=\sum_i x_i a_i,\qquad x_i\in A,
\]
would express $\ell$ as a $k\langle W\rangle$-linear combination of the
$Q^2$-coefficients of the $a_i$.  This would make
$L=k\langle W,W^{-1}\rangle$ a finite $k\langle W\rangle$-module, which is
impossible: a finite algebra is integral, whereas a monic equation for
$W^{-1}$, after multiplication by a power of $W$, would put
$1$ in the proper ideal $(W)\subset k\langle W\rangle$.  Thus $J$ is not
finitely generated, and $A$ is not noetherian.
\end{proof}


\section{A nonuniform rational localisation}

We next compute the rational chart which witnesses the failure of stable
uniformity.

Consider the datum
\[
  \alpha_{\mathrm{in}}=(W;\varpi).
\]
The ideal \((\varpi,W)\) is open, since \(\varpi\) is a unit in \(A\).  Let
\[
  A_{\mathrm{in}}=A\left\langle\frac{W}{\varpi}\right\rangle
\]
be the corresponding rational localisation on $|W|\leq |\varpi|$, and write
$X=W/\varpi$.  By the ring-of-definition construction, $A_{\mathrm{in}}$ is the separated completion of the subring generated by $A_0$ and $X$.

\begin{proposition}\label{prop:bad-chart}
There is a canonical isomorphism of complete topological rings
\[
  A_{\mathrm{in}}\cong k\langle X,Q\rangle/(Q^2),
  \qquad W\longmapsto\varpi X,\quad Q\longmapsto Q.
\]
\end{proposition}

\begin{proof}
Put
\[
  T=k\langle X,Q\rangle/(Q^2).
\]
Take the two-jet map $A\to B$ and then rescale its disc variable by
$W\mapsto\varpi X$.  This gives a continuous map
\[
  \theta:A\longrightarrow T,\qquad
  f_0(W)+Qf_1(W)+Q^2h\longmapsto
  f_0(\varpi X)+Qf_1(\varpi X).
\]
The element $\theta(\varpi)$ is a unit and
$\theta(W)/\theta(\varpi)=X$ is power-bounded.  The universal property first
extends $\theta$ to the algebraic rational localisation and then, by
continuity, to a map
\[
  \Phi:A_{\mathrm{in}}\longrightarrow T.
\]

We next record why the two-jet map loses precisely the terms divisible by
$Q^2$ after completion.  If the constant and linear $Q$-coefficients of
$y\in A$ vanish, then $y_n=W^{-n}y$ still belongs to $A$, and the $y_n$
have bounded norm.  In $A_{\mathrm{in}}$ one has
\[
  \overline y=\varpi^nX^n\overline{y_n}\qquad(n\geq0).
\]
The powers of $X$ are bounded and $\varpi^n\to0$, so $\overline y=0$.
In particular, the image $\overline Q$ of $Q$ satisfies
$\overline Q^{\,2}=0$.

Let $\overline X=\overline W/\varpi$.  Since $\overline X$ is power-bounded,
evaluation at $\overline X$ gives a continuous map
$k\langle X\rangle\to A_{\mathrm{in}}$.  We may therefore define
\[
  \Psi:T\longrightarrow A_{\mathrm{in}},\qquad
  f+Qg\longmapsto f(\overline X)+g(\overline X)\overline Q.
\]
To check the first composite, decompose
$a\in A$ as $a=f_0(W)+Qf_1(W)+y$, where the first two coefficients lie in
$k\langle W\rangle$ and the first two $Q$-coefficients of $y$ vanish.  The
preceding paragraph kills $y$, and hence
$\Psi(\theta(a))=\overline a$.  This proves
$\Psi\Phi=1$ on the dense algebraic rational localisation, and therefore on
$A_{\mathrm{in}}$.  In the other direction, $\Phi\Psi$ fixes constants, $X$
and $Q$.  It is consequently the identity first on polynomials, then on
$k\langle X\rangle$ by density, and finally on $T$.  Thus $\Phi$ and $\Psi$
are inverse continuous homomorphisms.
\end{proof}

\begin{proposition}\label{prop:not-stably-uniform}
The ring $A$ is not stably uniform.
\end{proposition}

\begin{proof}
In $T=k\langle X,Q\rangle/(Q^2)$, every element of the line $kQ$ is nilpotent and hence power-bounded.  This line is not bounded with respect to the ring of definition $T_0$: for every $r\geq 0$,
\[
  \varpi^r\bigl(\varpi^{-(r+1)}Q\bigr)=\varpi^{-1}Q\notin T_0.
\]
Thus $T$ is not uniform.  By \cref{prop:bad-chart}, a rational localisation of the
uniform ring $A$ is nonuniform.
\end{proof}


\section{Localisation of the Milnor square}

In this section we prove that rational localisation preserves the defining pullback.
We first record the strict Koszul estimate needed for the proof.

\subsection{Koszul complexes for rational localisations}

Let $E$ be one of $B,C,D$, with its coefficientwise Gauss norm and the
displayed unit ball $E_0$.  Let $f_1,\dots,f_m,g\in E_0$ generate the unit
ideal in $E$.  Put
\[
  T_E=E\langle T_1,\dots,T_m\rangle,
  \qquad
  T_{E,0}=E_0\langle T_1,\dots,T_m\rangle,
\]
and set $r_i=gT_i-f_i$.

We give $T_E$ the topology for which
$\{\varpi^nT_{E,0}\}_{n\geq0}$ is a basis of neighbourhoods of zero.  The term
$\bigwedge^qT_E^m$ has the finite-module topology defined by the lattice
$\bigwedge^qT_{E,0}^m$; equivalently, in the standard exterior basis it has the
finite product topology.  Images have the subspace topology.  Thus a
differential is strict when it is open onto its image, or equivalently when its
coimage is homeomorphic to its image.

Applied to $E\in\{B,C,D\}$, the strict Koszul regularity theorem of
Bambozzi--Kremnizer
\cite[Notation~4.3, Definitions~4.5--4.6, Proposition~4.9,
Lemma~4.11, and Corollary~4.13]{BK} gives the exactness, strictness, and
closed-image assertions below.  We include the direct proof followed in the
formalisation because the two estimates relative to the displayed lattices
are not stated there and are needed below.

\begin{lemma}\label{lem:koszul}
The Koszul complex
\[
  0\longrightarrow \bigwedge^mT_E^m\longrightarrow\cdots
  \longrightarrow\bigwedge^2T_E^m\xrightarrow{d_{2,E}}
  T_E^m\xrightarrow{d_{1,E}}T_E
\]
is exact in positive degrees.  Every differential is continuous and strict onto
its image, and every image is closed.  In particular, if
\[
  I_E=(r_1,\ldots,r_m)=\im(d_{1,E})\subset T_E,
  \qquad
  d_{1,E}(u_1,\ldots,u_m)=\sum_i u_ir_i,
\]
then $I_E$ is closed.  Moreover, there exists
$h_E\in\mathbb Z_{\geq0}$ such that
\begin{equation}\label{eq:ideal-denominators}
  \varpi^{h_E}(I_E\cap T_{E,0})
  \subset d_{1,E}(T_{E,0}^m).
\end{equation}
There also exists $z_E\in\mathbb Z_{\geq0}$ such that
\begin{equation}\label{eq:syzygy-denominators}
  \varpi^{z_E}\bigl(\ker(d_{1,E})\cap T_{E,0}^m\bigr)
  \subset d_{2,E}\bigl(\textstyle\bigwedge^2 T_{E,0}^m\bigr).
\end{equation}
\end{lemma}

\begin{proof}
Over the polynomial ring $E[T_1,\ldots,T_m]$, the sequence
$gT_i-f_i$ has Koszul complex exact in positive degrees.  The map
\[
  E[T_1,\ldots,T_m]\longrightarrow T_E
\]
is flat: this is the usual identification of $T_E$ with the
$\varpi$-adic completion of $E_0[T_1,\ldots,T_m]$, followed by inverting
$\varpi$.  Tensoring the polynomial Koszul complex with $T_E$, and using the
standard base-change identification of its finite free terms, proves
positive-degree exactness.  Notice that this does not assert exactness in
degree zero.

Each differential is continuous because its coordinates are finite sums of
multiplication maps.  The image of $d_{1,E}$ is closed by the noetherian
closed-ideal theorem.  In higher degree, exactness identifies an image with
the kernel of the next continuous differential, and hence that image is
closed.  The nonarchimedean open mapping theorem, applied to each
differential with its closed image, now gives a constant $C_q$ such that every
element of the image has a preimage of norm at most $C_q$ times its norm.
This proves strictness.

For $d_{1,E}$ choose $h_E$ so that
$|\varpi|^{h_E}C\leq1$, where $C$ is a controlled-lift constant for
$d_{1,E}$.  Scaling a controlled lift of an element of
$I_E\cap T_{E,0}$ gives \eqref{eq:ideal-denominators}.  Applying the same
argument to $d_{2,E}$, with its own constant and using
$\im(d_{2,E})=\ker(d_{1,E})$, gives
\eqref{eq:syzygy-denominators}.
\end{proof}


\subsection{The defining ideals under pullback}

Fix a rational datum
\[
  \alpha=(f_1,\dots,f_m;g)
\]
in $A$.  Multiplying all entries by a common power of $\varpi$, we assume
$f_1,\dots,f_m,g\in A_0$.  For $E\in\{A,B,C,D\}$ use the notation
$T_E$, $T_{E,0}$, and $d_{1,E}$ introduced above, and write
\[
  I_E=(gT_1-f_1,\ldots,gT_m-f_m)\subset T_E.
\]

The exact integral row \eqref{eq:integral-milnor} and the coefficientwise section of
$C_0\to D_0$ give an exact sequence
\begin{equation}\label{eq:tate-ambient}
  0\longrightarrow T_{A,0}\longrightarrow
  T_{B,0}\oplus T_{C,0}\longrightarrow T_{D,0}\longrightarrow 0.
\end{equation}
The map $T_{C,0}\to T_{D,0}$ is surjective, as are the induced maps on
finite free modules and their exterior powers.  Exactness is coefficientwise:
a compatible restricted coefficient pair lifts to $A_0$, and the lifted
coefficients still tend to zero $\varpi$-adically.

\begin{lemma}\label{lem:graph-pullback}
The ideal $I_A$ is closed in $T_A$, and the canonical map
\[
  I_A\longrightarrow I_B\times_{I_D}I_C
\]
is an isomorphism.  Moreover,
\begin{equation}\label{eq:ideal-row}
  0\longrightarrow I_A\longrightarrow I_B\oplus I_C
  \longrightarrow I_D\longrightarrow 0
\end{equation}
is strict exact.
\end{lemma}

\begin{proof}
We use the controlled lifts supplied in the proof of
\cref{lem:koszul}.  Let
$(x_B,x_C)\in I_B\times_{I_D}I_C$, and choose lifts
$u_B\in T_B^m$ and $u_C\in T_C^m$ of $x_B$ and $x_C$ whose norms are bounded
by fixed multiples of $\lVert x_B\rVert$ and $\lVert x_C\rVert$.  In $T_D^m$
the difference
\[
  w=(u_B)_D-(u_C)_D
\]
belongs to $\ker(d_{1,D})$.  Choose a controlled lift
$v_D\in\bigwedge^2T_D^m$ with $d_{2,D}(v_D)=w$, and lift $v_D$
coefficientwise to $v_C\in\bigwedge^2T_C^m$.  The coefficientwise section
$T_C\to T_D$ preserves the norm.  Replacing $u_C$ by
\[
  u_C'=u_C+d_{2,C}(v_C)
\]
does not change $d_{1,C}(u_C)$, and the images of $u_B$ and $u_C'$ in
$T_D^m$ now agree.  Exactness of \eqref{eq:tate-ambient}, applied
coefficientwise, gives $u_A\in T_A^m$ mapping to this pair.  Then
$d_{1,A}(u_A)$ maps to $(x_B,x_C)$.

The same construction keeps track of the norms.  The controlled lifts for
$d_{1,B}$, $d_{1,C}$ and $d_{2,D}$, continuity of $d_{2,C}$, and the
norm-one coefficientwise section give a constant $C$, independent of
$(x_B,x_C)$, for which the lift $x_A=d_{1,A}(u_A)$ may be chosen with
\[
  \lVert x_A\rVert\leq
  C\max\{\lVert x_B\rVert,\lVert x_C\rVert\}.
\]
The map $T_A\to T_B\oplus T_C$ is injective, so this proves both the
algebraic pullback assertion and the boundedness of its inverse.

It remains to see that $I_A$ is closed.  The ideals $I_B$ and $I_C$ are
closed by \cref{lem:koszul}.  If $x$ lies in the closure of $I_A$, its images
lie in $I_B$ and $I_C$.  The controlled pullback just proved gives
$x_A\in I_A$ with the same two images as $x$.  Injectivity of
$T_A\to T_B\oplus T_C$ gives $x=x_A$, as required.

Finally, let $y\in I_D$.  A controlled lift of $y$ through $d_{1,D}$,
followed by the coefficientwise section $T_D^m\to T_C^m$, produces
$x_C\in I_C$ mapping to $y$, with norm bounded by a fixed multiple of
$\lVert y\rVert$.  Thus $I_B\oplus I_C\to I_D$ is a strict surjection.
Its kernel is $I_A$ by the pullback assertion, and the preceding norm estimate
makes the inclusion on the left strict.  This proves
\eqref{eq:ideal-row}.
\end{proof}

\begin{proposition}\label{prop:localized-milnor}
For every rational datum $\alpha$ in $A$, the canonical sequence
\begin{equation}\label{eq:localized-row}
  0\longrightarrow A_\alpha\longrightarrow
  B_\alpha\oplus C_\alpha\longrightarrow D_\alpha\longrightarrow 0
\end{equation}
is strict exact, and $C_\alpha\to D_\alpha$ is a strict surjection.  Equivalently,
\[
  A_\alpha\cong B_\alpha\times_{D_\alpha}C_\alpha.
\]
The coefficient maps used in these identifications commute with rational
restriction.
\end{proposition}

\begin{proof}
The ideals $I_B,I_C,I_D$ are closed by \cref{lem:koszul}, and $I_A$ is
closed by \cref{lem:graph-pullback}.  Hence \eqref{eq:graph-model} gives
\[
  E_\alpha=T_E/I_E
  \qquad(E=A,B,C,D).
\]

We first prove algebraic exactness.  If a representative $p\in T_A$ maps
into both $I_B$ and $I_C$, the controlled pullback in
\cref{lem:graph-pullback} gives an element of $I_A$ with the same two images.
Injectivity of $T_A\to T_B\oplus T_C$ then shows that $p\in I_A$.  This gives
injectivity on the left.  Conversely, take compatible classes in
$T_B/I_B$ and $T_C/I_C$ and choose representatives $b$ and $c$.  Their
defect in $T_D$ belongs to $I_D$.  Lift this defect to an element of $I_C$
using the strict surjection in \eqref{eq:ideal-row}, and add it to $c$.
The corrected representatives agree in $T_D$, so
\eqref{eq:tate-ambient} lifts them to $T_A$.  This proves
\eqref{eq:localized-row}.

For the topology on the left, let $x\in A_\alpha$ and put
\[
  M=\max\{\lVert x_B\rVert,\lVert x_C\rVert\}.
\]
Choose representatives $b\in T_B$ and $c\in T_C$ whose norms are
arbitrarily close to the quotient norms of $x_B$ and $x_C$.  Their defect
lies in $I_D$.  The controlled surjection $I_C\to I_D$ gives a correction
$y\in I_C$ whose norm is bounded by a fixed multiple of the norm of this
defect.  Thus $b$ and $c+y$ agree in $T_D$ and lift to
$p\in T_A$.  The max-norm identity for the ambient pullback gives, after
letting the error in the representatives tend to zero,
\[
  \lVert x\rVert\leq C M
\]
for a constant $C$ depending only on the datum.  The two quotient maps are
norm-nonincreasing, so this estimate says precisely that
$A_\alpha\to B_\alpha\oplus C_\alpha$ is a topological embedding.

The coefficientwise section $T_D\to T_C$ shows directly that
$C_\alpha\to D_\alpha$ is surjective.  It also shows that
$T_C\to T_D$ is open; passing to quotient seminorms proves that the induced
map is open.  Together with continuity, this makes
$C_\alpha\to D_\alpha$ a strict surjection.

All maps used above come from the coefficient maps
$A\to B,C$ and $B,C\to D$.  More generally, a homomorphism $E\to E'$ carrying
a rational datum $\alpha$ to $\alpha'$ induces a coefficientwise map
$T_E\to T_{E'}$ which sends the relations for $\alpha$ to those for
$\alpha'$.  The pullback identifications therefore commute with these maps.
Applying this to rational restriction, and then using presentation independence
and transitivity, proves the asserted compatibility under refinement.
\end{proof}

\section{Sheafiness}

\begin{theorem}\label{thm:sheafy}
The Huber pair $(A,A^\circ)$ is sheafy.
\end{theorem}

\begin{proof}
Give $B,C,D$ their maximal rings of integral elements.  The max/Gauss estimates \eqref{eq:structural-map-bounds} show that every structural map is norm-nonincreasing.  Consequently it sends bounded subsets to bounded subsets; in particular, if $x$ is power-bounded, then the set of powers of its image is the image of the bounded set $\{x^n:n\geq0\}$.  The structural maps therefore send power-bounded elements to power-bounded elements and define morphisms of Huber pairs.

Each of $B,C,D$ is an affinoid $k$-algebra and is therefore sheafy by
Huber's theorem \cite[Theorem 2.2]{Huber94}.  Let
$U=\bigcup_iU_i$ in $\Spa(A,A^\circ)$.  Push the datum for $U$ and the data
for the $U_i$ to $B,C,D$.  If a section over $U$ restricts to zero on every
$U_i$, its images over $B$ and $C$ vanish by separatedness there.
Injectivity of the localised Milnor row from
\cref{prop:localized-milnor} then shows that the original section is zero.

Now take a matching family over the $U_i$.  Its images glue over $B$ and
$C$.  The two glued sections have the same image over $D$, since their
restrictions agree on the pushed cover and the presheaf for $D$ is
separated.  Exactness of the localised Milnor row over $U$ gives a section
over $A$, and injectivity of the corresponding rows over the $U_i$ shows
that it restricts to the given family.

It remains to retain the topology.  The image of the product restriction
map is the matching-family equaliser just described, which is closed in the
finite product of the section rings.  The restriction map is continuous and
injective, so the nonarchimedean closed-range theorem, using the
topologically nilpotent unit of $A$, makes it a topological embedding.  This
proves the sheaf condition on finite rational covers.  The
sheaf-on-a-rational-basis theorem \cite[Tag 009O]{Stacks} then gives the
all-open topological sheaf condition for $(A,A^\circ)$.
\end{proof}

\begin{proof}[Proof of \cref{thm:main}]
Completeness, uniformity, nonnoetherianity, and the domain property are
\cref{prop:uniform-domain}; sheafiness is
\cref{thm:sheafy}; and failure of stable uniformity is
\cref{prop:not-stably-uniform}.
\end{proof}


\section{A second example}\label{sec:second-example}

The finite-jet construction above loses uniformity because separated completion creates a
nilpotent.  We now give a second construction in which the bad rational localisation is an
integral domain.  The loss of uniformity is instead caused by an unbounded family of
nonnilpotent power-bounded elements.

\begin{theorem}\label{thm:second-example}
There is a complete Tate $k$-algebra $\mathcal A$ with the following properties.
\begin{enumerate}
\item $\mathcal A$ is a uniform, nonnoetherian integral domain and
      $\mathcal A^\circ=\mathcal A_0$ for the ring of definition displayed below.
\item The Huber pair $(\mathcal A,\mathcal A^\circ)$ is sheafy.
\item The rational localisation
      \[
        \mathcal B=\mathcal A\left\langle\frac{W}{\varpi}\right\rangle
      \]
      is an integral domain but is not uniform.
\end{enumerate}
Consequently $\mathcal A$ is not stably uniform.  In particular, failure of
stable uniformity need not be caused by a nilpotent in the bad rational
localisation.
\end{theorem}

\subsection{The weighted-parity algebra}

We write $\mathbb N=\{0,1,2,\ldots\}$ and put
$I=\mathbb N_{>0}$.  For a finite-support multi-index
$\nu=(\nu_n)_{n\in I}\in\mathbb N^{(I)}$, put
\begin{equation}\label{eq:parity-weight}
  \omega(\nu)=\sum_{\substack{n\geq1\\ \nu_n\text{ odd}}} n.
\end{equation}
Consider the submonoid
\begin{equation}\label{eq:parity-monoid}
  S=\left\{(a,\nu)\in\mathbb N\times\mathbb N^{(I)}:
       a\geq\omega(\nu)\right\}.
\end{equation}
The inequality
\begin{equation}\label{eq:omega-subadditive}
  \omega(\nu+\mu)\leq\omega(\nu)+\omega(\mu)
\end{equation}
shows that $S$ is closed under addition.  It is locally finite: a fixed element of $S$ has
only finitely many decompositions as a sum of two elements of $S$.

Let $k[S]$ be the monoid algebra with basis $W^aU^\nu$, where
$U^\nu=\prod_nU_n^{\nu_n}$.  Give it the Gauss norm
\begin{equation}\label{eq:parity-gauss-norm}
  \left\lVert\sum_{(a,\nu)\in S}c_{a,\nu}W^aU^\nu\right\rVert_G
  =\sup_{a,\nu}|c_{a,\nu}|.
\end{equation}
Its completion is
\begin{equation}\label{eq:parity-algebra}
  \mathcal A=
  \left\{
    \sum_{(a,\nu)\in S}c_{a,\nu}W^aU^\nu:
    c_{a,\nu}\in k,\quad c_{a,\nu}\longrightarrow0
  \right\},
\end{equation}
where convergence to zero is over the discrete index set $S$.  Local finiteness makes the
coefficientwise convolution finite, and the submultiplicativity of the Gauss norm on $k[S]$
extends multiplication continuously to $\mathcal A$.  Put
\begin{equation}\label{eq:parity-ring-of-definition}
  \mathcal A_0=
  \left\{
    \sum_{(a,\nu)\in S}c_{a,\nu}W^aU^\nu\in\mathcal A:
    c_{a,\nu}\in k^\circ
  \right\}.
\end{equation}
Then $\mathcal A_0$ is the closed unit ball for \eqref{eq:parity-gauss-norm}, it is
$\varpi$-adically complete, and
$\mathcal A=\mathcal A_0[1/\varpi]$.

Put
\[
  Y_n=W^nU_n,\qquad Z_n=U_n^2,
\]
and, for $N\in\mathbb N$, let $\mathcal A^{(N)}\subset\mathcal A$ be the
closed subring supported on monomials involving only
$W,U_1,\ldots,U_N$.  Thus
$\mathcal A^{(0)}=k\langle W\rangle$.  Let
$\mathcal T^{(N)}\subset\mathcal A^{(N)}$ be the subring in which every
$U_n$ has even exponent.

\begin{lemma}[Finite-variable subrings]\label{lem:finite-variable-subrings}
For every $N\in\mathbb N$, substitution $Z_n\mapsto U_n^2$ gives a
norm-preserving ring isomorphism
\[
  k\langle W,Z_1,\ldots,Z_N\rangle
  \xrightarrow{\ \sim\ }\mathcal T^{(N)}.
\]
For $N=0$ this is the identity on $k\langle W\rangle$.
Every $f\in\mathcal A^{(N)}$ can be written
\begin{equation}\label{eq:finite-variable-decomposition}
  f=\sum_{\epsilon\in\{0,1\}^N} f_\epsilon Y^\epsilon,
  \qquad f_\epsilon\in\mathcal T^{(N)},
  \qquad Y^\epsilon=\prod_{n=1}^NY_n^{\epsilon_n}.
\end{equation}
Consequently $\mathcal A^{(N)}$ is finite over $\mathcal T^{(N)}$ and is
a strongly noetherian Tate ring.  The inclusions
$\mathcal A^{(N)}\to\mathcal A^{(M)}\to\mathcal A$ for $N\leq M$
preserve the norm.  Moreover, for every $f\in\mathcal A$ and $\ell\geq0$
there are $N$ and $f_N\in\mathcal A^{(N)}$ such that
\[
  \lVert f-f_N\rVert_G\leq |\varpi|^\ell\lVert f\rVert_G.
\]
In particular,
\begin{equation}\label{eq:A-completion-finite-variables}
  \mathcal A=\widehat{\bigcup_{N\geq0}\mathcal A^{(N)}}.
\end{equation}
\end{lemma}

\begin{proof}
On even exponents, halving the $U_n$-coordinates identifies the support
monoid with that of the Tate algebra
$k\langle W,Z_1,\ldots,Z_N\rangle$; the coefficient map and its inverse
preserve the Gauss norm.  Write each exponent uniquely as
$\nu_n=2q_n+\epsilon_n$, with $\epsilon_n\in\{0,1\}$.  Since
$\omega(\nu)=\sum_{\epsilon_n=1}n$, every allowed monomial factors as
\[
  W^aU^\nu
  =W^{a-\omega(\nu)}
   \prod_{n=1}^NZ_n^{q_n}
   \prod_{n=1}^NY_n^{\epsilon_n}.
\]
Grouping the monomials with the same tuple
$(\epsilon_1,\ldots,\epsilon_N)$ gives
\eqref{eq:finite-variable-decomposition}.  The finitely many
$Y^\epsilon$ therefore generate $\mathcal A^{(N)}$ over
$\mathcal T^{(N)}$.
The same argument after adjoining any finite number of Tate variables gives
strong noetherianity.  The inclusions are coefficientwise and hence
isometric.  Finally, approximate $f$ to the required precision by a finite
sum of monomials and choose $N$ so that only
$U_1,\ldots,U_N$ occur in that sum.  This gives the stated estimate and
\eqref{eq:A-completion-finite-variables}.
\end{proof}

\begin{proposition}\label{prop:parity-uniform-domain}
The ring $\mathcal A$ is a complete uniform Tate $k$-algebra and an integral domain, with
\[
  \mathcal A^\circ=\mathcal A_0.
\]
\end{proposition}

\begin{proof}
The Gauss norm on the countable restricted Tate algebra
$k\langle W,U_1,U_2,\ldots\rangle$ is multiplicative.  To see this, note
first that the norm of a nonzero restricted series is attained.  For each of
two such series choose, among the norm-attaining exponents, the
lexicographically largest one.  In the coefficient of the sum of these two
exponents, the product of the chosen coefficients strictly dominates every
other convolution term.  The ultrametric inequality therefore gives
\[
  \lVert fg\rVert_G=\lVert f\rVert_G\lVert g\rVert_G.
\]
This is the countable restricted-series argument used in the formalisation,
whose underlying Lean implementation is adapted from work of William Coram
\cite{CoramXia}.  The support algebra $\mathcal A$ is a closed isometric
subring of this Tate algebra, so its Gauss norm is multiplicative and it is a
domain.

If $x\in\mathcal A_0$, all powers of $x$ remain in $\mathcal A_0$.  If
$x\notin\mathcal A_0$, multiplicativity gives
$\lVert x^m\rVert_G=\lVert x\rVert_G^m$, which is unbounded.  Hence
$\mathcal A^\circ=\mathcal A_0$.  Completeness and the Tate property follow from
\eqref{eq:parity-algebra} and the topologically nilpotent unit $\varpi$.
\end{proof}

\begin{proposition}\label{prop:parity-nonnoetherian}
The ring $\mathcal A$ is not noetherian.
\end{proposition}

\begin{proof}
For $m\geq1$, let
\[
  I_m=(Z_1,\ldots,Z_m)\subset\mathcal A.
\]
Define a norm-nonincreasing homomorphism
\[
  \psi_m:\mathcal A\longrightarrow k\langle T\rangle,
\]
by retaining only the coefficients of the monomials
$U_{m+1}^{2q}$ and sending that monomial to $T^q$.  The support condition
shows directly that this is multiplicative: a product can contribute such a
monomial only when both factors contribute monomials of the same form.
Thus $\psi_m$ sends $W$ and every $Y_j$ to zero, sends $Z_{m+1}$ to $T$,
and sends every other $Z_j$ to zero.  It kills $I_m$ but not $Z_{m+1}$.
Therefore
\[
  I_1\subsetneq I_2\subsetneq I_3\subsetneq\cdots,
\]
so $\mathcal A$ is not noetherian.
\end{proof}

\subsection{A reduced nonuniform rational chart}

The ideal $(W,\varpi)$ is open because $\varpi$ is a unit of $\mathcal A$.
Define the complete $k^\circ$-algebra
\begin{equation}\label{eq:weighted-chart-lattice}
  \mathcal B_0=
  \left\{
    \sum_{\substack{\nu\in\mathbb N^{(I)},\ d\geq\omega(\nu)}}
      b_{d,\nu}X^dU^\nu:
    b_{d,\nu}\in\varpi^{\omega(\nu)}k^\circ,
    \quad
    \varpi^{-\omega(\nu)}b_{d,\nu}\longrightarrow0
  \right\},
\end{equation}
with norm
\begin{equation}\label{eq:weighted-chart-norm}
  \left\lVert\sum b_{d,\nu}X^dU^\nu\right\rVert_{\mathrm{in}}
  =\sup_{d,\nu}|\varpi|^{-\omega(\nu)}|b_{d,\nu}|.
\end{equation}
The subadditivity \eqref{eq:omega-subadditive} shows that $\mathcal B_0$ is a ring.  Put
$\mathcal B=\mathcal B_0[1/\varpi]$.

\begin{proposition}\label{prop:weighted-chart-identification}
There is a bicontinuous ring isomorphism
\begin{equation}\label{eq:weighted-chart-identification}
  \mathcal A\left\langle\frac{W}{\varpi}\right\rangle
  \xrightarrow{\ \sim\ }\mathcal B,
  \qquad W\longmapsto\varpi X.
\end{equation}
\end{proposition}

\begin{proof}
The datum $(W;\varpi)$ lies in
$\mathcal A^{(0)}=k\langle W\rangle$.  Let
\[
  Q=
  \mathcal A^{(0)}\langle T\rangle/
  \overline{(\varpi T-W)}
\]
be its separated completed quotient.  Evaluation at $W=\varpi X$ and $T=X$
gives a continuous map $Q\to k\langle X\rangle$.  Conversely, evaluation of
$k\langle X\rangle$ at the image of $T$ gives a continuous map in the other
direction.  The two maps are inverse by polynomial density, and the norm
estimates in the two universal properties show that the isomorphism is
isometric.

Apply the coefficientwise-localisation construction of
\cref{prop:coefficientwise-localization} with $N=0$.  It identifies the
rational chart with the $c_0$-sum over $Q$ obtained by grouping terms
according to their $U$-multi-index.  After replacing $Q$ by
$k\langle X\rangle$, this is the model
\eqref{eq:weighted-chart-lattice}.  Its sup norm is
\eqref{eq:weighted-chart-norm}: the term indexed by $\nu$ is multiplied by
$\varpi^{\omega(\nu)}$.  Evaluation on
$\mathcal A^{(0)}$ sends $W$ to $\varpi X$.
\end{proof}

\begin{proposition}\label{prop:weighted-chart-domain-nonuniform}
The ring $\mathcal B$ is an integral domain but is not uniform.
\end{proposition}

\begin{proof}
Group a series in $\mathcal B$ by its $U$-multi-index.  This gives a
coefficientwise homomorphism
\[
  \mathcal B\longrightarrow
  k\langle X\rangle[[U_1,U_2,\ldots]].
\]
It is injective, and the target is a domain because
$k\langle X\rangle$ is a domain.  Hence $\mathcal B$ is a domain.  In the
coefficientwise description below, this is the embedding of the
corresponding ring of coefficient families that tend to zero into formal
multivariable power series; the factor $\varpi X$ appearing in the
multiplication is a non-zero-divisor.

Inside $\mathcal A$ write $Y_n=W^nU_n$ and put
\[
  T_n=\varpi^{-n}Y_n=X^nU_n\in\mathcal B.
\]
Formula \eqref{eq:weighted-chart-norm} gives
\[
  \lVert T_n\rVert_{\mathrm{in}}=|\varpi|^{-n},
  \qquad
  \lVert T_n^2\rVert_{\mathrm{in}}=1.
\]
It follows by separating even and odd powers that the powers of each fixed
$T_n$ are bounded.  Thus every $T_n$ is power-bounded.  The family
$(T_n)_{n\geq1}$ is not bounded, since its norms
$|\varpi|^{-n}$ are unbounded.  Hence
$\mathcal B^\circ$ is unbounded and $\mathcal B$ is not uniform.
\end{proof}

\subsection{Small perturbations of rational data}

\begin{lemma}[Small perturbation lemma]\label{lem:small-perturbation}
Let $(E,E^+)$ be a complete Tate Huber pair and let
$\alpha=(f_1,\ldots,f_m;g)$ be a rational datum in $E$.  There is a
neighbourhood $V$ of zero in $E$ such that, whenever
\[
  g'-g,\ f_1'-f_1,\ldots,f_m'-f_m\in V,
\]
the tuple $\alpha'=(f_1',\ldots,f_m';g')$ is a rational datum and
$\alpha$ and $\alpha'$ define the same rational subset of
$\Spa(E,E^+)$.  Their completed rational localisations are canonically
bicontinuously isomorphic over $E$.
\end{lemma}

\begin{proof}
Apply \cite[Lemma~3.10]{Huber93} after including $g$ among the
numerators; see also \cite[Remark~2.4.7]{KL}.
The assertion about completed rational localisations follows from their
universal property \cite[Proposition~1.3]{Huber94}.
\end{proof}

\subsection{Rational localisations defined over finitely many variables}

Fix $N\in\mathbb N$.  Let
$\mu=(\mu_n)_{n>N}$ be a finitely supported family of nonnegative integers
and put
\begin{equation}\label{eq:coefficient-monomial}
  e_\mu=W^{\omega(\mu)}U^\mu,
\end{equation}
where $\mu$ is extended by zero for $n\leq N$.  Every
$f\in\mathcal A$ has a unique convergent expression
\begin{equation}\label{eq:coefficient-decomposition}
  f=\sum_\mu f_\mu e_\mu,
  \qquad
  f_\mu\in\mathcal A^{(N)},
  \qquad
  \lVert f_\mu\rVert_G\longrightarrow0,
\end{equation}
and
\[
  \lVert f\rVert_G=\sup_\mu\lVert f_\mu\rVert_G.
\]
Multiplication is determined by
\begin{equation}\label{eq:coefficient-multiplication}
  e_\mu e_\lambda
  =W^{\omega(\mu)+\omega(\lambda)-\omega(\mu+\lambda)}
   e_{\mu+\lambda}.
\end{equation}
The map
\begin{equation}\label{eq:finite-variable-retraction}
  \rho_N:\mathcal A\longrightarrow\mathcal A^{(N)},
  \qquad
  \sum_\mu f_\mu e_\mu\longmapsto f_0,
\end{equation}
is a norm-nonincreasing ring retraction.

\begin{proposition}[Coefficientwise localisation]\label{prop:coefficientwise-localization}
Let $\alpha$ be a rational datum whose entries lie in
$\mathcal A^{(N)}$, and put
$P=(\mathcal A^{(N)})_\alpha$.  There is a bicontinuous ring isomorphism
\begin{equation}\label{eq:coefficientwise-localization}
  \mathcal A_\alpha
  \xrightarrow{\ \sim\ }
  \left\{
    \sum_\mu p_\mu e_\mu:
    p_\mu\in P,\quad \lVert p_\mu\rVert\longrightarrow0
  \right\}.
\end{equation}
The ring on the right has the supremum norm, and its multiplication is
given by \eqref{eq:coefficient-multiplication}, with $W$ replaced by its
image in $P$.  If a rational refinement is also represented by data in
$\mathcal A^{(N)}$ and $P\to Q$ is the corresponding restriction map,
then the induced restriction sends
\begin{equation}\label{eq:coefficientwise-restriction}
  \sum_\mu p_\mu e_\mu
  \longmapsto
  \sum_\mu (p_\mu|_Q)e_\mu.
\end{equation}
\end{proposition}

\begin{proof}
Let
\[
  R=\mathcal A^{(N)}\langle T_1,\ldots,T_m\rangle,
  \qquad
  J=\overline{(gT_i-f_i)_i}\subset R,
  \qquad \widetilde P=R/J.
\]
Evaluation at $T_i=f_i/g$ identifies $\widetilde P$ with $P$.

The expression \eqref{eq:coefficient-decomposition} continues to hold after
adjoining the variables $T_i$.  Applying the quotient map
$R\to\widetilde P$ to
each coefficient gives a continuous homomorphism
\[
  \Phi:\mathcal A_\alpha\longrightarrow
  \left\{
    \sum_\mu q_\mu e_\mu:
    q_\mu\in\widetilde P,\quad \lVert q_\mu\rVert\longrightarrow0
  \right\}.
\]
Conversely, the canonical map from $\widetilde P$ to
$\mathcal A_\alpha$ may be applied
to every coefficient.  If
$\lVert q_\mu\rVert\to0$, the resulting series
\[
  \sum_\mu q_\mu e_\mu
\]
converges in $\mathcal A_\alpha$, and defines a continuous homomorphism
$\Psi$ in the other direction.  The multiplication rule
\eqref{eq:coefficient-multiplication} shows that both maps are ring
homomorphisms.

The two composites are the identity on sums with finitely many nonzero
coefficients, each represented by a polynomial.  Such sums are dense on
both sides, so continuity gives
$\Psi\Phi=1$ and $\Phi\Psi=1$.  Composing with
$\widetilde P\cong P$ proves
\eqref{eq:coefficientwise-localization}.

For a rational refinement represented in $\mathcal A^{(N)}$, the two
evaluation maps commute with the restriction map $P\to Q$.  The preceding
construction therefore gives \eqref{eq:coefficientwise-restriction}.
\end{proof}

\begin{corollary}[Finite-variable presentation]\label{cor:finite-variable-presentation}
Every rational localisation of $\mathcal A$ is bicontinuously isomorphic
to a ring of the form
\begin{equation}\label{eq:arbitrary-finite-variable-presentation}
  \left\{
    \sum_\mu p_\mu e_\mu:
    p_\mu\in P,\quad \lVert p_\mu\rVert\longrightarrow0
  \right\},
\end{equation}
where $P$ is a rational localisation of one of the strongly noetherian Tate
rings $\mathcal A^{(N)}$.
\end{corollary}

\begin{proof}
Start with a rational datum
$\alpha=(f_1,\ldots,f_m;g)$ in $\mathcal A$.  Multiplying all entries by
one power of $\varpi$ does not change the rational subset or localisation,
so assume that $g,f_1,\ldots,f_m$ lie in $\mathcal A_0$.  Choose a
neighbourhood $V$ as in \cref{lem:small-perturbation}.  By
\eqref{eq:A-completion-finite-variables}, for some $N$ there are
$g',f_1',\ldots,f_m'\in\mathcal A^{(N)}\cap\mathcal A_0$ such that
\[
  g'-g,\ f_1'-f_1,\ldots,f_m'-f_m\in V.
\]
By \cref{lem:small-perturbation}, the datum
$\alpha'=(f_1',\ldots,f_m';g')$ gives the same rational subset and the same
completed localisation.  It is rational in $\mathcal A$, so applying
$\rho_N$ to a Bezout identity for its entries shows that it is already a
rational datum in $\mathcal A^{(N)}$.  Now apply
\cref{prop:coefficientwise-localization}.
\end{proof}

\subsection{Sheafiness}

\begin{theorem}\label{thm:parity-sheafy}
The Huber pair $(\mathcal A,\mathcal A^\circ)$ is sheafy.
\end{theorem}

\begin{proof}
We first work with a finite rational cover $U=\bigcup_iU_i$.  Choose one
$N$ large enough to approximate the datum for $U$ and the data for all the
$U_i$ in $\mathcal A^{(N)}$.  The small perturbation lemma does not change
the rational domains, and \cref{prop:coefficientwise-localization} gives
compatible identifications
\[
  \mathcal O_X(U)\cong
  \left\{\sum_\mu p_\mu e_\mu:
    p_\mu\in P,\ \lVert p_\mu\rVert\to0\right\},
\]
\[
  \mathcal O_X(U_i)\cong
  \left\{\sum_\mu p_{i,\mu}e_\mu:
    p_{i,\mu}\in P_i,\ \lVert p_{i,\mu}\rVert\to0\right\},
\]
with coefficientwise restriction maps.  Here $P$ is a rational localisation
of the strongly noetherian Tate ring $\mathcal A^{(N)}$, and the $P_i$ are
rational localisations of $P$.  The domains defined by the $P_i$ cover
$\Spa(P,P^\circ)$.  Indeed, compose a point of $\Spa(P,P^\circ)$ with the
coefficient projection from the ring in the first display to $P$.  The
resulting point lies in one of the original $U_i$, and restricting back to
$P$ gives the required point of the corresponding domain.

Let $(x_i)_i$ be a matching family and write
$x_i=\sum_\mu x_{i,\mu}e_\mu$.  Since $P$ is sheafy,
the matching family $(x_{i,\mu})_i$ glues to a unique $q_\mu\in P$ for every
$\mu$.  It remains to check that $\lVert q_\mu\rVert\to0$.  The restrictions
of the $q_\mu$ to every $P_i$ tend to zero because the coefficients of each
$x_i$ do.  The product restriction map
\[
  P\longrightarrow\prod_iP_i
\]
is a topological embedding, so it reflects convergence to zero.  Hence
$q_\mu\to0$, and $\sum_\mu q_\mu e_\mu$ is the required section over $U$.
The same coefficientwise argument proves separation.

The image of the product restriction map for $\mathcal A$ is therefore the
matching-family equaliser.  This equaliser is closed, and the restriction map
is continuous and injective.  The nonarchimedean closed-range theorem over
the Tate ring $\mathcal A$ now makes it a topological embedding.  We have
proved the topological sheaf condition on finite rational covers.  The
general rational-basis theorem upgrades this to the all-open structure
presheaf, and its plus-ring independence gives the same conclusion for every
ring of integral elements.  In particular, $(\mathcal A,\mathcal A^\circ)$
is sheafy.
\end{proof}

\begin{proof}[Proof of \cref{thm:second-example}]
The uniform domain assertion is \cref{prop:parity-uniform-domain}, and nonnoetherianity is
\cref{prop:parity-nonnoetherian}.  Sheafiness is
\cref{thm:parity-sheafy}.  Finally,
\cref{prop:weighted-chart-identification,prop:weighted-chart-domain-nonuniform} identify the
rational chart $|W|\leq|\varpi|$ as a nonuniform integral domain.
\end{proof}

\begin{remark}[Relation with the first construction]\label{rem:second-example-relation}
The finite-jet and weighted-parity examples use different mechanisms.  In the finite-jet ring,
all terms of $Q$-degree at least two become infinitely
$\varpi$-divisible, and the bad chart acquires a nilpotent.  In the
weighted-parity ring, the data defining any finite rational calculation can
be approximated inside some strongly noetherian
$\mathcal A^{(N)}$, after which localisation and restriction are computed
coefficient by coefficient in the remaining variables.  Its bad chart remains a domain, and
nonuniformity is witnessed by the unbounded family $(T_n)_{n\geq1}$ of nonnilpotent
power-bounded elements.  The use of an infinite monomial support and of unbounded
power-bounded directions is in the tradition of the examples of Buzzard--Verberkmoes and
Mihara; the coefficientwise calculation over the rings
$\mathcal A^{(N)}$ is the additional feature that makes sheafiness
accessible here.
\end{remark}

\section{How the examples were found}\label{sec:discovery}

The author had for some time been interested in finding a counterexample to
Problem~7 of the \emph{Nonarchimedean Scottish Book}.  Our attempts had
centred on the example of
\cite[Section~4.5, Proposition~17]{BV}, which is uniform but not stably
uniform.  We believe that this example is in fact sheafy, but we have not
found a proof.

In January 2026 we began to experiment with AI tools, initially ChatGPT 5.2
Pro.  Buzzard--Verberkmoes reduce exactness for a two-member Laurent cover to
a boundedness condition on the relevant rings of definition
\cite[Lemma~2]{BV}.  For a uniform ring this follows from boundedness of the
power-bounded elements \cite[Lemma~3 and Corollary~4]{BV}.  Their reduction
from general rational covers to Laurent covers is given in
\cite[Lemma~8 and the proof of Theorem~7]{BV}.  For sheafiness these checks
must also be made after rational localisation.  Their Theorem~7 cannot be
applied directly to Proposition~17, since one of the rational localisations
in that example is not uniform.  Our plan was instead to ask ChatGPT to
write down many explicit rational covers and check the required exactness
directly for each one, in the hope that a pattern, and eventually a proof,
would emerge.

After a couple of weeks ChatGPT 5.2 Pro claimed to have a complete proof.
The author then began the slow process of understanding this proof, only to
discover, with the help of ChatGPT 5.3, that one of its arguments used the
following assertion:

\emph{If $\phi:A\to B$ is a continuous homomorphism of Hausdorff
topological abelian groups and $A$ is complete, then $\phi(A)$ is closed in
$B$.}

This assertion is false, and its use completely broke the purported proof.

Motivated by this, we began in parallel an AI-assisted Lean formalisation of
the parts of the theory of adic spaces needed here, including the theorem
that strongly noetherian Tate rings are sheafy.  The aim was a practical
one: future AI-generated proofs could be formalised before the human author
had spent days understanding an argument containing a fatal error near its
end.

We then returned to writing down covers and looking for a pattern.  This
continued through ChatGPT 5.3 and subsequent versions.  When ChatGPT 5.6 Sol
became available, we asked, almost in passing, whether it could prove the
result we wanted or \emph{find a new example}.  It produced the examples
studied here.  By then the Lean development was far enough advanced that we
could formalise the arguments quickly and check that they were correct.



\appendix
\section{Lean formalisation and verification}\label{sec:formalisation}

This appendix records the Lean formalisation of the two examples.  The
finite-jet declarations are taken from AINTLIB commit
\texttt{b007a4f3d}, and the weighted-parity declarations from commit
\texttt{090a28921} \cite{AINTLIB}.  For each example Lean proves the
Tate structure, uniformity, domain and nonnoetherianity statements,
sheafiness, and failure of stable uniformity, together with the rational
chart calculation used in the proof.

Neither pinned snapshot was available from the public AINTLIB repository at
the time of writing.  The interactive version of the paper includes each
declaration used below, extracted from the appropriate commit.  The
\href{https://github.com/CBirkbeck/uniform-sheafy-tate-domains/blob/main/PAPERFORGE_CHECK.md}{verification report}
records the Lean and mathlib versions, the build, and the axiom checks.

\subsection{Sheafiness}

The finite-jet proof first establishes a condition on finite rational covers.
Its historical internal class is \texttt{IsSheafy}, in the
\texttt{ValuationSpectrum} namespace and defined in
\href{https://cbirkbeck.github.io/uniform-sheafy-tate-domains/lean/AdicSpaces/declarations/ValuationSpectrum.IsSheafy.html}{\texttt{StructureSheaf.lean}}.
Here
\href{https://cbirkbeck.github.io/uniform-sheafy-tate-domains/lean/AdicSpaces/declarations/ValuationSpectrum.RationalCoveringData.html}{\texttt{RationalCoveringData A}}
consists of a rational datum for the base and finitely many data whose rational
subsets cover it.  The separate condition
\href{https://cbirkbeck.github.io/uniform-sheafy-tate-domains/lean/AdicSpaces/declarations/ValuationSpectrum.RationalCoveringData.IsRational.html}{\texttt{C.IsRational}}
says that the numerator set in the base and in each member of the cover
generates an open ideal.  For a Tate ring this is equivalent to generating the
unit ideal.  Finally,
\href{https://cbirkbeck.github.io/uniform-sheafy-tate-domains/lean/AdicSpaces/declarations/ValuationSpectrum.presheafValue.html}{\texttt{presheafValue D}}
is the completed rational localisation $A\langle T/s\rangle$ attached to $D$.

\begin{lstlisting}
class IsSheafy (A : Type u) [CommRing A] [TopologicalSpace A]
    [IsTopologicalRing A] [inst₁ : PlusSubring A]
    [inst₂ : IsHuberRing A] [T2Space A] [NonarchimedeanRing A]
    [letI : UniformSpace A :=
      IsTopologicalAddGroup.rightUniformSpace A; CompleteSpace A]
    [IsRingOfIntegralElements (A⁺)] : Prop where
  embedding : ∀ (C : RationalCoveringData A), C.IsRational →
    Topology.IsEmbedding (productRestrictionSub A C)
  gluing : ∀ (C : RationalCoveringData A), C.IsRational →
    ∀ (f : ∀ (D : ↥C.covers), presheafValue D.1),
    (∀ (D₁ D₂ : ↥C.covers) (D₃ : RationalLocData A)
      (h₃₁ : rationalOpen D₃.T D₃.s ⊆
        rationalOpen D₁.1.T D₁.1.s)
      (h₃₂ : rationalOpen D₃.T D₃.s ⊆
        rationalOpen D₂.1.T D₂.1.s),
      restrictionMap D₁.1 D₃ h₃₁ (f D₁) =
        restrictionMap D₂.1 D₃ h₃₂ (f D₂)) →
    ∃ x : presheafValue C.base, ∀ (D : ↥C.covers),
      restrictionMap C.base D.1 (C.hsubset D.1 D.2) x = f D
\end{lstlisting}

The two fields say that restriction to a finite rational cover is a
topological embedding and that every matching family glues.  Compatibility is
written using common rational refinements; the exact-intersection theorem
identifies this with the usual compatibility on pairwise intersections.  Thus
\texttt{IsSheafy A} is the topological sheaf condition on the rational basis.

For an arbitrary open subset $V\subseteq\Spa(A,A^+)$,
\href{https://cbirkbeck.github.io/uniform-sheafy-tate-domains/lean/AdicSpaces/declarations/ValuationSpectrum.limitSections.html}{\texttt{limitSections V}}
is the ring of compatible sections over the rational subdomains contained in
$V$, or equivalently their projective limit.  If $W\subseteq V$, then
\href{https://cbirkbeck.github.io/uniform-sheafy-tate-domains/lean/AdicSpaces/declarations/ValuationSpectrum.limitRestrict.html}{\texttt{limitRestrict}}
is the map $\mathcal O_X(V)\to\mathcal O_X(W)$ obtained by reindexing this
family.  The corresponding all-open condition is \texttt{IsLimitSheaf}, from
\href{https://cbirkbeck.github.io/uniform-sheafy-tate-domains/lean/AdicSpaces/declarations/ValuationSpectrum.IsLimitSheaf.html}{\texttt{SheafyPair.lean}}.

The class
\href{https://cbirkbeck.github.io/uniform-sheafy-tate-domains/lean/AdicSpaces/declarations/ValuationSpectrum.HasLocLiftPowerBounded.html}{\texttt{HasLocLiftPowerBounded}}
is used internally to construct the restriction maps between completed
rational localisations.  For an inclusion $R(D')\subseteq R(D)$, it says that
the denominator of $D$ becomes a unit on $R(D')$ and that the fractions
$t/s$, for $t$ in the numerator set of $D$, are power-bounded there.  This
class is supplied automatically for the complete Huber rings considered here;
it is not an additional hypothesis in our results.

\begin{lstlisting}
structure IsLimitSheaf (A : Type u) [CommRing A] [TopologicalSpace A]
    [IsTopologicalRing A] [PlusSubring A] [IsHuberRing A]
    [HasLocLiftPowerBounded A] : Prop where
  injective : ∀ {V : Opens ↥(Spa A A⁺)} {ι : Type u}
    {U : ι → Opens ↥(Spa A A⁺)}
    (hle : ∀ i, U i ≤ V)
    (_ : (V : Set ↥(Spa A A⁺)) ⊆ ⋃ i, (U i : Set ↥(Spa A A⁺)))
    {x y : ↥(limitSections V)},
    (∀ i, limitRestrict (hle i) x = limitRestrict (hle i) y) → x = y
  glue : ∀ {V : Opens ↥(Spa A A⁺)} {ι : Type u}
    {U : ι → Opens ↥(Spa A A⁺)}
    (hle : ∀ i, U i ≤ V)
    (_ : (V : Set ↥(Spa A A⁺)) ⊆ ⋃ i, (U i : Set ↥(Spa A A⁺)))
    (s : ∀ i, ↥(limitSections (U i))),
    (∀ i j,
      limitRestrict (inf_le_left (a := U i) (b := U j)) (s i) =
      limitRestrict (inf_le_right (a := U i) (b := U j)) (s j)) →
    ∃ x : ↥(limitSections V), ∀ i, limitRestrict (hle i) x = s i
  isEmbedding : ∀ {V : Opens ↥(Spa A A⁺)} {ι : Type u}
    {U : ι → Opens ↥(Spa A A⁺)}
    (hle : ∀ i, U i ≤ V)
    (_ : (V : Set ↥(Spa A A⁺)) ⊆ ⋃ i, (U i : Set ↥(Spa A A⁺))),
    Topology.IsEmbedding (limitRestrictProd hle)
\end{lstlisting}

The fields \texttt{injective} and \texttt{glue} are the usual sheaf axioms on
arbitrary opens.  The final field says that the resulting map to the
matching-family equaliser is also a topological embedding.  This is the extra
condition needed for a sheaf of complete topological rings.

For complete Tate pairs the rational-basis and all-open predicates are
equivalent:

\begin{lstlisting}
theorem isSheafy_iff_isLimitSheaf :
    IsSheafy A ↔ IsLimitSheaf A :=
  ⟨fun _ => isLimitSheaf_of_isSheafy,
    isSheafy_of_isLimitSheaf⟩
\end{lstlisting}

The forward implication is the sheaf-on-a-rational-basis theorem, and the
converse is restriction to rational covers.  The
\href{https://cbirkbeck.github.io/uniform-sheafy-tate-domains/lean/AdicSpaces/declarations/ValuationSpectrum.isSheafOfTopologicalRings_iff_isLimitSheaf.html}{structure-presheaf equivalence}
identifies \texttt{IsLimitSheaf A} with
Wedhorn's topological-ring formulation \cite[Remark~8.20]{Wedhorn}.
The development also proves Mathlib's categorical sheaf predicate for its
\texttt{CompleteTopCommRingCat}-valued presheaf.  This predicate does record
the topology, since its morphisms are continuous; Wedhorn's formulation
instead tests against all topological rings, rather than only complete
separated ones.

For a complete Tate ring the following definitions fix a ring of integral
elements and then remove the choices of plus ring and completion.  They are in
\href{https://cbirkbeck.github.io/uniform-sheafy-tate-domains/lean/AdicSpaces/declarations/ValuationSpectrum.IsSheafyFor.html}{\texttt{SheafyRing.lean}}.

\begin{lstlisting}
def IsSheafyFor (Aplus : RingOfIntegralElements A) : Prop :=
  letI := Aplus.toPlusSubring
  haveI : IsRingOfIntegralElements (A⁺ : Subring A) := Aplus.2
  haveI : HasLocLiftPowerBounded A := hasLocLiftPowerBounded_faithful
  IsLimitSheaf A

def IsSheafyComplete : Prop :=
  ∀ Aplus : RingOfIntegralElements A, IsSheafyFor A Aplus

def IsSheafyTateRing : Prop :=
  ∀ (P : PairOfDefinition A)
    (Bplus : RingOfIntegralElements (CompletionModel A P)),
    IsSheafyFor (CompletionModel A P) Bplus
\end{lstlisting}

The predicate \texttt{IsSheafyFor A Aplus} says that the Huber pair
$(A,A^+)$ is sheafy.  For complete Tate rings, the
\href{https://cbirkbeck.github.io/uniform-sheafy-tate-domains/lean/AdicSpaces/declarations/ValuationSpectrum.isSheafyFor_iff_isSheafyComplete.html}{chosen-plus-ring
equivalence} shows that this is independent of the chosen ring of integral
elements.  The equivalences in
\href{https://cbirkbeck.github.io/uniform-sheafy-tate-domains/lean/AdicSpaces/declarations/ValuationSpectrum.isSheafyTateRing_iff_for_completionModel.html}{\texttt{SheafyCompletionModel.lean}}
remove the choice of completion.  Consequently \texttt{IsSheafyTateRing} is
the analytic specialisation of Wedhorn's Definition~8.26
\cite[Definition~8.26]{Wedhorn}.

\subsection{Strong noetherianity}

The definitions below are from
\href{https://cbirkbeck.github.io/uniform-sheafy-tate-domains/lean/AdicSpaces/declarations/IsStronglyNoetherian.html}{\texttt{RestrictedPowerSeries.lean}}.

\begin{lstlisting}
def MvPowerSeries.IsRestricted {k : ℕ} {A : Type*}
    [CommRing A] [TopologicalSpace A]
    (f : MvPowerSeries (Fin k) A) : Prop :=
  Tendsto (fun s : Fin k →₀ ℕ => MvPowerSeries.coeff s f)
    cofinite (nhds 0)

class IsStronglyNoetherian (A : Type*) [CommRing A]
    [TopologicalSpace A] [NonarchimedeanRing A] : Prop where
  isNoetherianRing_restricted : ∀ k : ℕ,
    IsNoetherianRing (restrictedMvPowerSeriesSubring k A)
\end{lstlisting}

The cofinite convergence condition is the usual condition that the
coefficients tend to zero, so
\texttt{restrictedMvPowerSeriesSubring k A} represents
$A\langle T_1,\ldots,T_k\rangle$.  The second declaration therefore says that
these rings are noetherian for every $k$.  Since $A$ is already complete, this
is Wedhorn's definition of strong noetherianity
\cite[Proposition and Definition~6.36]{Wedhorn}.

\subsection{The finite-jet example}

We finally record the definition of the example itself.  In the following
code, \texttt{PowerSeries.Restricted R 1} is the radius-one Tate algebra in one
variable over $R$, \texttt{RestrictedLaurent K} is
$K\langle W,W^{-1}\rangle$, and \texttt{DualNumber R} is $R[Q]/(Q^2)$.
The source is
\href{https://cbirkbeck.github.io/uniform-sheafy-tate-domains/lean/AdicSpaces/declarations/FiniteJetOver.JetA.html}{\texttt{FJP/Over/JetRings.lean}}.

\begin{lstlisting}
abbrev L : Type u := RestrictedLaurent K

abbrev JetC : Type u := PowerSeries.Restricted (L K) (1 : ℝ)

abbrev JetB : Type u :=
  DualNumber (PowerSeries.Restricted K (1 : ℝ))

abbrev JetD : Type u := DualNumber (L K)

noncomputable def jetSupport : Subring (JetC K) where
  carrier := {f |
    qCoeff K 0 f ∈ nonnegSubring K ∧
    qCoeff K 1 f ∈ nonnegSubring K}
  zero_mem' := by
    constructor <;> rw [qCoeff_zero] <;> exact zero_mem _
  one_mem' := by
    constructor <;> rw [qCoeff_one]
    · rw [if_pos rfl]; exact one_mem _
    · rw [if_neg one_ne_zero]; exact zero_mem _
  add_mem' := fun {f g} hf hg => by
    constructor <;> rw [qCoeff_add]
    · exact add_mem hf.1 hg.1
    · exact add_mem hf.2 hg.2
  neg_mem' := fun {f} hf => by
    constructor <;> rw [qCoeff_neg]
    · exact neg_mem hf.1
    · exact neg_mem hf.2
  mul_mem' := fun {f g} hf hg => by
    constructor
    · rw [qCoeff_zero_mul]
      exact mul_mem hf.1 hg.1
    · rw [qCoeff_one_mul]
      exact add_mem (mul_mem hf.1 hg.2) (mul_mem hf.2 hg.1)

abbrev JetA : Type u := ↥(jetSupport K)
\end{lstlisting}

The predicate \texttt{nonnegSubring K} cuts out
$K\langle W\rangle$ inside $K\langle W,W^{-1}\rangle$.  Consequently the
carrier condition says that the coefficients of $Q^0$ and $Q^1$ have no
negative powers of $W$, while the coefficients of $Q^n$ for $n\geq2$ are
unrestricted.  Thus \texttt{JetA K} is
\[
  \left\{f_0(W)+Qf_1(W)+Q^2h(W,W^{-1},Q)\ \middle|\
    \substack{f_0,f_1\in K\langle W\rangle\\
              h\in L\langle Q\rangle}\right\},
\]
which is the coefficient description of $A$ in
\cref{eq:coefficient-A0} after inverting $\varpi$.

The file also constructs the four maps in the Milnor square and proves
\texttt{milnorRow\_exact}: every compatible pair in
\texttt{JetB K} and \texttt{JetC K} comes from a unique element of
\texttt{JetA K}.  This identifies \texttt{JetA K} with
$B\times_D C$.  The sheaf argument produces the finite-rational-cover
condition displayed above:

\begin{lstlisting}
include ϖ hK₀ in
theorem isSheafy_JetA :
    ValuationSpectrum.IsSheafy (JetA K) where
  embedding := fun C hC =>
    productRestrictionSub_isEmbedding_JetA ϖ hK₀ C hC
  gluing := fun C hC f hcompat =>
    gluing_JetA ϖ hK₀ C hC f hcompat
\end{lstlisting}

Here $\varpi$ is a \texttt{Uniformizer K}, and the remaining input is the
noetherianity of \texttt{unitBall K}.  This is the version of sheafiness proved
directly for the example.  The equivalences above give the all-open and
choice-independent statements.  The main theorems are collected in
\href{https://cbirkbeck.github.io/uniform-sheafy-tate-domains/lean/AdicSpaces/declarations/FiniteJetOver.finiteJet_isSheafyFor_of_dvr.html}{\texttt{SheafyEndpoints}},
and include the following.

\begin{lstlisting}
theorem finiteJet_isDomain : IsDomain (JetA K) :=
  inferInstance

theorem finiteJet_not_noetherian :
    ¬ IsNoetherianRing (JetA K) :=
  not_isNoetherianRing_JetA K

theorem finiteJet_isUniform_of_dvr
    [IsDiscreteValuationRing 𝒪[K]] :
    TopologicalRing.IsUniform (JetA K) :=
  finiteJet_isUniform K (Uniformizer.ofDVR K)

theorem finiteJet_not_stablyUniform_of_dvr
    [IsDiscreteValuationRing 𝒪[K]] :
    ¬ TopologicalRing.IsStablyUniform (JetA K) :=
  finiteJet_not_stablyUniform K (Uniformizer.ofDVR K)

theorem finiteJet_isSheafyFor_of_dvr
    [IsDiscreteValuationRing 𝒪[K]] :
    IsSheafyFor (JetA K) (finiteJetPlus K) :=
  finiteJet_isSheafyFor K (Uniformizer.ofDVR K)
    (isNoetherianRing_unitBall K)

theorem finiteJet_isSheafyTateRing_of_dvr
    [IsDiscreteValuationRing 𝒪[K]] :
    IsSheafyTateRing (JetA K) :=
  finiteJet_isSheafyTateRing K (Uniformizer.ofDVR K)
    (isNoetherianRing_unitBall K)
\end{lstlisting}

Here $K$ is any complete nontrivially normed ultrametric field for which
$\mathcal O[K]=\{x:\lVert x\rVert\leq1\}$ is a discrete valuation ring.  No
local compactness, finiteness of the residue field, characteristic, or
perfection assumption is used.  The theorem
\texttt{finiteJet\_isSheafyFor\_of\_dvr} proves that $(A,A^\circ)$ is sheafy,
whilst \texttt{finiteJet\_isSheafyTateRing\_of\_dvr} removes the choices of
completion and ring of integral elements.  The corresponding all-open
statement is
\[
\begin{split}
&\texttt{TopCat.Presheaf.IsSheafOfTopologicalRings}\\[-2pt]
&\qquad
\texttt{(ValuationSpectrum.structurePresheaf (JetA K))},
\end{split}
\]
The
\href{https://cbirkbeck.github.io/uniform-sheafy-tate-domains/lean/AdicSpaces/declarations/FiniteJetOver.finiteJet_structurePresheaf_isSheafOfTopologicalRings_all_of_dvr.html}{corresponding Lean theorem}
proves this statement directly.

\subsection{Scope of the finite-jet formalisation}

The paper's $A_0$ is represented by the Gauss-norm unit ball.
The theorem \texttt{isPowerBounded\_JetA\_iff} identifies this unit ball with
the power-bounded elements, giving $A^\circ=A_0$.  Completeness and the Tate
structure are supplied by the corresponding instances.  As in the
weighted-parity development, the constant-series map from $K$ is explicit
and isometric but is not installed as an \texttt{Algebra K} instance.

For the bad chart, \texttt{chartEquiv} is an exported ring equivalence.
The two continuity theorems make it a homeomorphism, and the canonical-map
theorem identifies $W$ and $Q$.  It is not bundled as a
\texttt{K}-\texttt{AlgEquiv}, which is why \cref{prop:bad-chart} is stated for
complete topological rings.

The conclusions of \cref{lem:koszul} are formalised in the normed setting used
for the example: $E$ is one of $B,C,D$, $T_E$ has its Gauss norm,
$T_{E,0}$ is its unit ball, and the finite free terms have their sup norms.
Positive-degree exactness is transported from the polynomial sequence
by flat base change.  Closedness in positive degree comes from exactness and
continuity, whilst the degree-zero ideal is closed by the noetherian
closed-ideal theorem.  The closed-range theorem then gives controlled lifts
and strictness; scaling those lifts gives the two displayed denominator
inclusions.

The proof of \cref{lem:graph-pullback} in Lean uses these controlled lifts
directly.  Its localisation proof obtains the topological embedding from a
quotient-norm estimate, and proves openness of $C_\alpha\to D_\alpha$ from
the norm-preserving coefficientwise section.  The components of
\cref{prop:localized-milnor} are therefore separate declarations.  Finally,
the development proves the concrete finite-jet transfer
\texttt{isSheafy\_JetA}: it pushes finite rational covers to $B,C,D$, glues
there, and uses the localised Milnor row and a closed-range theorem.

The Koszul argument, the rational-chart calculation, and the transfer of
sheafiness through the concrete Milnor square are contained respectively in
\href{https://cbirkbeck.github.io/uniform-sheafy-tate-domains/lean/AdicSpaces/declarations/FiniteJet.GraphKoszul.koszulGraph_restricted_exact.html}{the
restricted-exactness file}
and
\href{https://cbirkbeck.github.io/uniform-sheafy-tate-domains/lean/AdicSpaces/declarations/FiniteJet.GraphKoszul.koszulDifferential_isStrict.html}{the
strictness and closed-image file},
\href{https://cbirkbeck.github.io/uniform-sheafy-tate-domains/lean/AdicSpaces/declarations/FiniteJetOver.chartEquiv.html}{\texttt{FJP/Over/Chart.lean}},
and
\href{https://cbirkbeck.github.io/uniform-sheafy-tate-domains/lean/AdicSpaces/declarations/FiniteJetOver.isSheafy_JetA.html}{\texttt{FJP/Over/SheafTransfer.lean}}.

\subsection{The weighted-parity example}

The second development uses the same restricted-series infrastructure.  Its
ambient ring is the countable restricted Tate algebra
$K\langle W,U_1,U_2,\ldots\rangle$.  The definitions
\href{https://cbirkbeck.github.io/uniform-sheafy-tate-domains/lean/AdicSpaces/declarations/WeightedParity.wpWeight.html}{\texttt{wpWeight}},
\href{https://cbirkbeck.github.io/uniform-sheafy-tate-domains/lean/AdicSpaces/declarations/WeightedParity.WPMem.html}{\texttt{WPMem}}, and
\href{https://cbirkbeck.github.io/uniform-sheafy-tate-domains/lean/AdicSpaces/declarations/WeightedParity.wpSupport.html}{\texttt{wpSupport}}
are the direct translation of \eqref{eq:parity-weight}--\eqref{eq:parity-algebra}.
The relevant part of the definition is as follows.

\begin{lstlisting}
def wpWeight (w : ℕ → ℕ) (t : ℕ →₀ ℕ) : ℕ :=
  ∑ n ∈ t.support,
    if (t n).mod 2 = 1 ∧ n ≠ 0 then w n else 0

def WPMem (w : ℕ → ℕ) (t : ℕ →₀ ℕ) : Prop :=
  wpWeight w t ≤ t 0

abbrev WPA : Type _ := ↥(wpSupport K w)

def idWeight : ℕ → ℕ := fun n => n

abbrev WPAid : Type _ := WPA K idWeight
\end{lstlisting}

The coordinate $t(0)$ is the exponent of $W$; the other coordinates are the
exponents of the $U_n$.  Thus \texttt{WPMem idWeight t} says exactly that
$a\geq\omega(\nu)$, and \texttt{wpSupport} requires the coefficients outside
this monoid to vanish.  It follows that
\href{https://cbirkbeck.github.io/uniform-sheafy-tate-domains/lean/AdicSpaces/declarations/WeightedParity.WPAid.html}{\texttt{WPAid K}}
is the ring $\mathcal A$ of \eqref{eq:parity-algebra}.  The maps
\href{https://cbirkbeck.github.io/uniform-sheafy-tate-domains/lean/AdicSpaces/declarations/WeightedParity.constA.html}{\texttt{constA}} and
\href{https://cbirkbeck.github.io/uniform-sheafy-tate-domains/lean/AdicSpaces/declarations/WeightedParity.norm_constA.html}{\texttt{norm\_constA}}
give the isometric embedding of $K$ as the constant series.  This map is
explicit in the development, although it is not bundled as an
\texttt{Algebra K} instance.

Suppose that $K$ is complete and nontrivially normed and that
$\mathcal O[K]$ is a DVR.  The principal exported statements are:

\begin{lstlisting}
theorem weightedParity_isUniform_of_dvr :
    IsUniform (WPAid K) :=
  weightedParity_isUniform (Uniformizer.ofDVR K)

theorem weightedParity_isDomain :
    IsDomain (WPAid K) :=
  inferInstance

theorem weightedParity_not_noetherian :
    ¬ IsNoetherianRing (WPAid K) :=
  not_isNoetherianRing_WPA K idWeight idWeight_pos

theorem weightedParity_powerBounded_eq_unitBall
    (ϖ : Uniformizer K) :
    powerBoundedSubring (WPAid K) =
      (FiniteJet.unitBall (WPAid K) : Set (WPAid K)) :=
  powerBoundedSubring_eq_unitBall ϖ

theorem weightedParity_isSheafyComplete_of_dvr :
    IsSheafyComplete (WPAid K) :=
  weightedParity_isSheafyComplete (Uniformizer.ofDVR K)
    (FiniteJetOver.isNoetherianRing_unitBall K)

theorem weightedParity_chart_isDomain
    (ϖ : Uniformizer K)
    (hK₀ : IsNoetherianRing (FiniteJet.unitBall K)) :
    IsDomain (presheafValue
      (chartDatum (w := idWeight) ϖ)) :=
  isDomain_chart ϖ hK₀

theorem weightedParity_chart_not_isUniform
    (ϖ : Uniformizer K)
    (hK₀ : IsNoetherianRing (FiniteJet.unitBall K)) :
    ¬ IsUniform (presheafValue
      (chartDatum (w := idWeight) ϖ)) :=
  not_isUniform_chart idWeight_unbounded ϖ hK₀

theorem weightedParity_not_stablyUniform_of_dvr :
    ¬ IsStablyUniform (WPAid K) :=
  weightedParity_not_stablyUniform (Uniformizer.ofDVR K)
    (FiniteJetOver.isNoetherianRing_unitBall K)
\end{lstlisting}

Here \texttt{IsSheafyComplete} is the predicate defined above: it quantifies
over every ring of integral elements.  In particular it proves the sheaf
condition for $(\mathcal A,\mathcal A^\circ)$.  The
\href{https://cbirkbeck.github.io/uniform-sheafy-tate-domains/lean/AdicSpaces/declarations/WeightedParity.wp_structurePresheaf_isSheaf_all.html}{all-open theorem}
also records the corresponding statement directly for the structure
presheaf.  The datum
\href{https://cbirkbeck.github.io/uniform-sheafy-tate-domains/lean/AdicSpaces/declarations/WeightedParity.chartDatum.html}{\texttt{chartDatum}}
is $(W;\varpi)$, and
\href{https://cbirkbeck.github.io/uniform-sheafy-tate-domains/lean/AdicSpaces/declarations/WeightedParity.chartDatum_isRational.html}{\texttt{chartDatum\_isRational}}
checks that it defines the chart $|W|\leq|\varpi|$.  The two chart theorems
then say precisely that its section ring is a domain but is not uniform.

The proof routes are the ones used in Section~7.  Multiplicativity of the
countable Gauss norm gives the domain and unit-ball statements.  For rational
data in $\mathcal A^{(N)}$, localisation is proved by taking the completed
quotient of $\mathcal A^{(N)}$, applying its evaluation map to each
coefficient in \eqref{eq:coefficient-decomposition}, and constructing the
inverse by summing coefficients that tend to zero; it does not invoke the
Koszul lemma.  For sheafiness, one $N$ is chosen for the base datum and all
members of a finite cover.  Strong noetherianity of
$\mathcal A^{(N)}$ glues each coefficient, and the restriction embedding
shows that the resulting coefficients tend to zero.
The final topological embedding is obtained from the closed matching
equaliser by the Tate closed-range theorem.  For the bad chart, the quotient
formed from $\mathcal A^{(0)}=K\langle W\rangle$ is
$K\langle X\rangle$, and the coefficientwise ring embeds in formal
multivariable power series over this domain.  The formalisation also
contains, under \texttt{WeightedParity.PerturbSetup}, a quantitative
norm-unit-ball version of the standard perturbation result cited in
\cref{lem:small-perturbation}.

These declarations, including the
\href{https://cbirkbeck.github.io/uniform-sheafy-tate-domains/lean/AdicSpaces/declarations/WeightedParity.weightedParity_not_stablyUniform_of_dvr.html}{final failure of stable uniformity},
have been checked to use only the usual logical axioms
\texttt{propext}, \texttt{Classical.choice}, and \texttt{Quot.sound}.

\subsection{Mathematical and code provenance}

The formal structure presheaf, rational-cover criterion, and the
strongly-noetherian sheaf theorem follow Huber's construction
\cite[Theorem~2.2 and Lemmas~2.3--2.6]{Huber94} and Wedhorn's systematic
formulation
\cite[Proposition and Definition~6.36, Remark~8.20, Definition~8.26,
Theorem~8.28(b), Corollary~8.32, and Lemmas~8.33--8.34]{Wedhorn}.
The passage between a fixed plus ring and all plus rings follows the standard
rational-refinement argument recorded by Kedlaya
\cite[Lemma~1.6.8 and Remark~1.6.9]{KedlayaAWS}.
Ben-Bassat--Kremnizer prove the strict one-variable resolutions used in their
treatment of rational localisations \cite[Lemmas~5.13--5.14]{BBK};
Bambozzi--Kremnizer formulate the general construction in terms of Koszul
complexes and prove strict Koszul regularity
\cite[Notation~4.3, Definitions~4.5--4.6, and
Corollary~4.13]{BK}.  The formalisation does not invoke either result.
Positive-degree exactness is instead proved by flat base change from the
polynomial sequence; closedness and the nonarchimedean open mapping theorem
\cite[Lemma~3.1]{KST} give controlled lifts.  Scaling those lifts yields the
bounded-denominator criterion used by Buzzard--Verberkmoes
\cite[Lemma~2]{BV}.

The development is written in Lean~4 \cite{Lean4} on top of mathlib
\cite{mathlib}.  The formalisation-specific files live in AINTLIB
\cite{AINTLIB}.  Its restricted power-series and restricted Laurent-series
infrastructure includes Apache-2.0-licensed code adapted from William Coram's
\texttt{PhD} repository, with equivalence code attributed there to Bingyu
Xia \cite{CoramXia}.  The vendored source headers and commit history retain
those attributions.  The finite-jet and weighted-parity definitions,
localisation proofs, and sheaf endpoints discussed in this appendix are the
AINTLIB-specific part of the development.

\begin{thebibliography}{99}

\bibitem{AINTLIB}
C.~Birkbeck and contributors,
\emph{AINTLIB: AdicSpaces formalisation},
Lean~4 source code, finite-jet snapshot
\texttt{b007a4f3d4226f00a684b402715aa542e2f0bcdc} and weighted-parity
snapshot
\texttt{090a289211deb69117413e329325fe819aa7dbc2}, 2026,
\url{https://github.com/CBirkbeck/AINTLIB}.

\bibitem{BBK}
O.~Ben-Bassat and K.~Kremnizer,
\emph{Non-Archimedean analytic geometry as relative algebraic geometry},
Ann. Fac. Sci. Toulouse Math. (6) \textbf{26} (2017), no.~1, 49--126,
\href{https://doi.org/10.5802/afst.1526}{doi:10.5802/afst.1526}.

\bibitem{BK}
F.~Bambozzi and K.~Kremnizer,
\emph{On the sheafyness property of spectra of Banach rings},
J. Lond. Math. Soc. (2) \textbf{109} (2024), no.~1, e12855,
\href{https://doi.org/10.1112/jlms.12855}{doi:10.1112/jlms.12855}.

\bibitem{BV}
K.~Buzzard and A.~Verberkmoes,
\emph{Stably uniform affinoids are sheafy},
J. Reine Angew. Math. \textbf{740} (2018), 25--39.

\bibitem{CoramXia}
W.~Coram and B.~Xia,
\emph{Restricted power-series and power-series equivalence code},
Lean source files in the \texttt{WilliamCoram/PhD} repository,
commit \texttt{e8fcf8fbff848a95475ab62ae2568cbb73961de8}, 2026,
\url{https://github.com/WilliamCoram/PhD}.

\bibitem{HK}
D.~Hansen and K.~S. Kedlaya,
\emph{Sheafiness criteria for Huber rings},
preprint, version of 23 April 2025,
\url{https://kskedlaya.org/papers/criteria.pdf}.

\bibitem{Huber93}
R.~Huber,
\emph{Continuous valuations},
Math. Z. \textbf{212} (1993), no.~3, 455--477,
\href{https://doi.org/10.1007/BF02571668}{doi:10.1007/BF02571668}.

\bibitem{Huber94}
R.~Huber,
\emph{A generalization of formal schemes and rigid analytic varieties},
Math. Z. \textbf{217} (1994), 513--551,
\href{https://doi.org/10.1007/BF02571959}{doi:10.1007/BF02571959}.


\bibitem{KedlayaAWS}
K.~S. Kedlaya,
\emph{Sheaves, stacks, and shtukas},
Arizona Winter School notes, 2017,
\url{https://swc-math.github.io/aws/2017/2017KedlayaNotes.pdf}.

\bibitem{KL}
K.~S. Kedlaya and R.~Liu,
\emph{Relative $p$-adic Hodge theory: Foundations},
Ast\'erisque \textbf{371} (2015),
\href{https://arxiv.org/abs/1301.0792}{arXiv:1301.0792}.

\bibitem{KST}
M.~Kerz, S.~Saito, and G.~Tamme,
\emph{$K$-theory of non-archimedean rings II},
Nagoya Math. J. \textbf{251} (2023), 669--685.

\bibitem{Lean4}
L.~de Moura and S.~Ullrich,
\emph{The Lean 4 theorem prover and programming language},
in \emph{Automated Deduction---CADE 28},
Lecture Notes in Comput. Sci., vol.~12699, Springer, 2021, 625--635,
\href{https://doi.org/10.1007/978-3-030-79876-5_37}{doi:10.1007/978-3-030-79876-5\_37}.

\bibitem{mathlib}
The mathlib Community,
\emph{The Lean mathematical library},
in \emph{Proceedings of the 9th ACM SIGPLAN International Conference on
Certified Programs and Proofs}, ACM, 2020, 367--381,
\href{https://doi.org/10.1145/3372885.3373824}{doi:10.1145/3372885.3373824}.

\bibitem{Milnor}
J.~Milnor,
\emph{Introduction to Algebraic $K$-Theory},
Annals of Mathematics Studies, vol.~72,
Princeton University Press, Princeton, 1971.

\bibitem{Mihara}
T.~Mihara,
\emph{On Tate's acyclicity and uniformity of Berkovich spectra and adic spectra},
Israel J. Math. \textbf{216} (2016), no.~1, 61--105.

\bibitem{Scottish}
K.~S. Kedlaya,
\emph{The Nonarchimedean Scottish Book}, Problem~7,
\href{https://scripts.mit.edu/~kedlaya/wiki/index.php?title=The_Nonarchimedean_Scottish_Book}{available online}.

\bibitem{Stacks}
The Stacks Project Authors,
\emph{The Stacks Project},
\url{https://stacks.math.columbia.edu}.

\bibitem{Wedhorn}
T.~Wedhorn,
\emph{Adic spaces},
2019,
\href{https://arxiv.org/abs/1910.05934}{arXiv:1910.05934}.

\end{thebibliography}

\end{document}
